\subsection{Emotion Recognition from Speech}
Speech emotion recognition has been extensively studied using both hand-crafted acoustic features \cite{eyben2010opensmile} and deep learning approaches \cite{chen2022wavlm,pepino2021emotion}. While WavLM \cite{chen2022wavlm} and other self-supervised models achieve state-of-the-art results on emotion classification, their applicability to generative tasks remains unexplored. Our work shows that for emotion-conditioned text generation, interpretable acoustic features may be preferable.

\subsection{Multimodal Large Language Models}
Recent work on multimodal LLMs has primarily focused on vision-language integration \cite{driess2023palm,alayrac2022flamingo,li2023blip}. PaLM-E \cite{driess2023palm} and Flamingo \cite{alayrac2022flamingo} demonstrate that grounding LLMs in visual perception improves reasoning and generation. However, audio modality, particularly emotional prosody, remains underexplored in LLM research.

\subsection{Knowledge Distillation for LLMs}
Knowledge distillation \cite{hinton2015distilling} has proven effective for transferring capabilities from large teacher models to smaller student models \cite{sanh2019distilbert}. Our work extends this paradigm to multimodal settings, where a teacher LLM's text responses are used to train a student model conditioned on both text and audio emotion features.

\subsection{Cross-Speaker Generalization}
Speaker-invariant feature learning is well-studied in speech recognition \cite{saon2013speaker} and speaker verification \cite{snyder2018x}, but its application to emotion recognition remains challenging. Prior work shows that speaker normalization \cite{zen2013statistical} and multi-speaker training \cite{dehak2011front} improve robustness, but our results suggest these approaches are insufficient for emotion-conditioned generation tasks.
